\begin{proof}[Proof of Theorem~\ref{thm:arith}]
Induction on the structure of the IR term $L(q)$, written $\hat{q}$ below. For
$\hat{q}\in\mathcal{K}$ this is Theorem~\ref{thm:kernel}, and for
$\hat{q}=\mathrm{Number}(c)$ both sides denote $c$.

$\mathrm{Scale}(c,e)$ emits $p_e * c$. Scalar--vector arithmetic is pointwise in both
languages and preserves the domain, so the claim follows from the hypothesis on $e$; the
same argument covers a $\mathrm{BinOp}$ one of whose operands is a $\mathrm{Number}$,
where $\chi$ takes the other operand's value.

The induction carries a second clause, which the steps below use and an earlier
version of this proof did not state: for every $\hat{e}\in\mathcal{A}$ the emitter
translates, the PromQL result of $p_e$ carries exactly the labels $\chi(\hat{e})$, and
its elements have pairwise distinct $\chi(\hat{e})$ values. It holds at each case for a
different reason. On $\mathcal{K}$ the emitted expression is an aggregation by
$\overline{k}$, so Lemma~\ref{lem:pq-matching}'s last clause gives the labels and the
grouping gives distinctness. $\mathrm{Number}$ is a scalar and carries none.
$\mathrm{Scale}$ is scalar--vector arithmetic, which preserves its operand's labels. And
the two $\mathrm{BinOp}$ cases below preserve it, the first because an intersection on
identical label sets does and the second because the result is keyed by
$\overline{k_l}$, which the left operand already satisfies.

Stating it matters because both steps below justified the clause by citing
Lemma~\ref{lem:pq-matching}'s last clause directly, and that clause is about
\emph{aggregations}: an operand of a $\mathrm{BinOp}$ in $\mathcal{A}$ need not be one,
since it may itself be a $\mathrm{Scale}$ or a $\mathrm{BinOp}$. The citation was to a
lemma that does not reach the operands it was invoked on.

For $\mathrm{BinOp}(\oplus,l,r)$ with both operands vectors, write
$\overline{k_l}=\chi(l)$ and $\overline{k_r}=\chi(r)$. By induction the two operand
results agree between the systems, so it is enough that the two languages pair them the
same way and the emitter asks for the pairing that does.

If $\overline{k_l}=\overline{k_r}=\overline{k}$, the emitter emits $p_l \oplus p_r$ with
no modifier. Each operand's PromQL result carries exactly the labels $\overline{k}$, by
Lemma~\ref{lem:pq-matching}'s last clause, so matching on identical label sets
\emph{is} matching on $\overline{k}$; the domain is the intersection and the values are
pointwise. Lemma~\ref{lem:dd-matching} gives Datadog the same intersection and the same
values.

If $\overline{k_r}\subsetneq\overline{k_l}$, the emitter emits
$p_l \oplus \texttt{on}(\overline{k_r})\ \texttt{group\_left()}\ p_r$. By
Lemma~\ref{lem:pq-matching} the result is keyed by $\overline{k_l}$, each element paired
with the unique element of $p_r$ agreeing on $\overline{k_r}$ --- unique by the
induction's second clause, which gives $p_r$'s elements distinct $\chi(r)$ values. The
earlier version said ``unique because $p_r$ is an aggregation by $\overline{k_r}$'',
which is false as written whenever $r$ is a $\mathrm{BinOp}$ or a $\mathrm{Scale}$. That is term for term the first clause of Lemma~\ref{lem:dd-matching}. The
symmetric case emits \texttt{group\_right()}. Taking $\overline{k_r}=\emptyset$ gives the
ungrouped operand, where \texttt{on()} matches everything and the broadcast is Datadog's.

If neither key set contains the other, $\chi(q)$ is undefined, the emitter refuses and
$q\notin\mathrm{dom}(T)$. Finally $C$ gains nothing, because the emitter attaches no
caveat at a $\mathrm{BinOp}$ or a $\mathrm{Scale}$: the modifiers it chooses are decided
by $\chi$, which is a function of the query text alone.

One case this does not cover, and the theorem's statement of exactness should be read
against it: a division whose right operand is zero. Both systems produce a
non-real reading there rather than a number, and which one --- an infinity, a
not-a-number, or an absent point --- is not compared anywhere in this paper.
Definition~\ref{def:d} is explicit that it is undefined on a non-real reading and that
every exactness test scores such a reading as a disagreement, so the theorem's equality
is an equality of numbers at instants where both sides produce numbers. Where they do
not, the claim is not that the two agree; it is that nothing here establishes they do.
\end{proof}
